% Page 100 - Labor Value Equations

\section*{Labor Value Calculations in Input-Output Framework}

\subsection*{Purchaser Price Equations:}
\[
M_p = (M_p)_p + (M_p)_t
\]

\[
TP^* = GO_p + GO_t
\]

Where:
\begin{itemize}
    \item $M_p$: Purchaser price of productive inputs
    \item $(M_p)_p$: Producer price component
    \item $(M_p)_t$: Trading margin component
    \item $TP^*$: Total product at purchaser prices
    \item $GO_p$: Gross output of productive sectors
    \item $GO_t$: Gross output of trading sector (trading margins)
\end{itemize}

\subsection*{Labor Value Calculation:}
\[
\lambda_j = hp_j + \sum_i \lambda_i \cdot app_{ij}
\]

Where:
\begin{itemize}
    \item $\lambda_j$ = labor value per unit output
    \item $hp_j$ = hours of productive labor per unit output
    \item $app_{ij}$ = quantity of the ith production input used per unit output
    \item $app_{ij} = \frac{[(M_p)_p]_{ij}}{X_j}$
    \item $X_j$ = quantity of output
\end{itemize}

\subsection*{Notes on Labor Value Adjustment:}
Ideally, adjust labor time flows for skill differences. If sectoral wage-rate differences are correlated with skill differences, wage rates can be used as first approximation. However, this can cause problems (see Wolff 1975, 1977 discussion in Section 6.1.2).